\begin{abstract}
В умовах плоскої деформації досліджено ініціювання відгалуження
V-подібної міжфазної тріщини зсуву у кутовій точці ламаної межі поділу
двох пружних ізотропних матеріалів. Тріщина розглядається як заданий
стаціонарний дефект із двома симетричними гілками, береги яких контактують без
тертя. Сингулярне поле в її кутовій точці використано як зовнішнє поле
для малої розтягувальної зони передруйнування, орієнтованої вздовж
бісектриси одного з матеріалів. Встановлено умови, за яких одночасно
забезпечуються стискальний контакт берегів міжфазної тріщини та
розтягувальне нормальне поле в матеріалі, де формується зона. За допомогою
перетворення Мелліна, зрощування асимптотик, методу Вінера--Гопфа та
умови регулярності напружень у вершині зони отримано замкнені аналітичні
співвідношення для її рівноважної довжини, розкриття, роботи сил опору
відриву та пов'язаного із зоною енергетичного параметра. Досліджено
залежність цих характеристик від кута зламу межі поділу, пружного
контрасту та рівня навантаження. Показано, що наявність зони
передруйнування послаблює сингулярність стаціонарної конфігурації з
V-подібною міжфазною тріщиною, хоча не усуває її повністю. За фіксованої
геометрії енергетична та деформаційна характеристики зони однозначно
пов'язані, тому відповідні критерії ініціювання відгалуження є
еквівалентними за узгодженого калібрування критичних параметрів.
Після незалежного задання відповідної критичної характеристики
матеріалу отримані результати дають змогу визначити критичне зовнішнє
навантаження ініціювання відгалуженої тріщини.
\end{abstract}

\noindent\textbf{Ключові слова:}
кусково-однорідне тіло; ламана межа поділу; V-подібна міжфазна тріщина
зсуву; зона передруйнування; ініціювання відгалуження; метод
Вінера--Гопфа.
